%% =============================================================================
%% sm_f3_diagnostics.tex — SM-F, Subsection F.3: Supplementary Diagnostics
%% Body content only (\subsection level); parent structure in sm_f.tex
%% =============================================================================

\subsection{Supplementary diagnostic figures and design effect analysis}
\label{smf:diagnostics}

This section presents three supplementary figures that visualize
bias, interval width, and design effects across all simulation
scenarios, and a full table of design effect ratios (DER) for
all 11 fixed-effect parameters.

\paragraph{Relative bias distributions.}
\Cref{smf:fig-bias} shows the distribution of relative bias
across parameters and scenarios.  E-UW exhibits the largest
absolute bias (especially for $\tau_{\mathrm{ext}}$ and
$\tau_{\mathrm{int}}$ under S4), while E-WT/E-WS improve
point estimation for fixed effects by incorporating survey
weights.  For the poverty slopes $\alpha_{\mathrm{pov}}$ and
$\beta_{\mathrm{pov}}$, the relative bias under E-WS remains
below 15\% across all scenarios.

\begin{figure}[htbp]
  \centering
  \includegraphics[width=0.95\textwidth]{SF7_simulation_bias.pdf}
  \caption{Relative bias (RB, \%) by parameter, estimator, and
    scenario.  Each panel corresponds to one of the five target
    parameters.  Dot positions indicate relative bias; color
    distinguishes scenarios (S0, S3, S4).  The vertical gray
    line marks zero bias.}
  \label{smf:fig-bias}
\end{figure}

\paragraph{Width ratio distributions.}
\Cref{smf:fig-wr} displays the distribution of width ratios
(WR = median CI width under E-WS / median CI width under E-WT)
across replications.  The progressive inflation from S0 through
S4 is clearly visible for all fixed-effect parameters.
The width ratio for $\log\kappa$ is similar to or slightly larger
than those for the poverty slopes, reflecting the nonlinear
dispersion score (Appendix~B).  Hyperparameter
width ratios are identically~$1.000$ (not shown) because no
sandwich correction is applied.

\begin{figure}[htbp]
  \centering
  \includegraphics[width=0.95\textwidth]{SF8_width_ratio_distribution.pdf}
  \caption{Distribution of width ratios across $R = 200$
    replications, by parameter and scenario.  Boxplots are overlaid
    with violin plots to show the full distributional shape.  The
    horizontal dashed line at WR = 1.0 marks the no-correction
    reference.}
  \label{smf:fig-wr}
\end{figure}

\paragraph{Design effect ratios.}
\Cref{smf:fig-der} summarizes the simulation DER across all 11
fixed-effect parameters (5 extensive, 5 intensive, and
$\log\kappa$) and three scenarios.  The DER monotonically
increases with the informativeness parameter
$\rho_{\mathrm{inc}}$, consistent with the theoretical prediction
in~\cref{smc:prop-der-classification}.  Urban density parameters
exhibit the largest DER in all scenarios (mean 3.25 in S0 up to
6.43 in S4), mirroring the empirical DER of~4.18 for
$\alpha_{\mathrm{urb}}$ in the NSECE case study
(\cref{smd:tab-der}).

\begin{figure}[htbp]
  \centering
  \includegraphics[width=0.95\textwidth]{SF9_DER_summary.pdf}
  \caption{Design effect ratios (DER) by parameter and scenario.
    Each dot represents the mean DER across 200 replications;
    color distinguishes scenarios.  Parameters are grouped by
    margin (extensive vs.\ intensive) with $\log\kappa$ as a
    separate category.}
  \label{smf:fig-der}
\end{figure}

\paragraph{Complete DER table.}
The table below reports the full distribution of DER across
replications for all 11 fixed-effect parameters and three
scenarios.  The mean, median, standard deviation, and range
(min--max) characterize the replication-to-replication
variability.  Two patterns are noteworthy: (i)~the DER standard
deviation increases with the scenario severity, indicating
greater sampling variability under heavier design effects; and
(ii)~the urban parameter consistently shows the widest DER
range, reflecting the spatial concentration of urban providers
that amplifies the design effect.

{\small
% latex table generated in R 4.5.1 by xtable 1.8-4 package
% Mon Feb 23 11:13:23 2026
\begin{table}[htbp]
\centering
\begin{tabular}{llrlllllll}
  \toprule
Scenario & Parameter & R & Mean & Median & SD & Min & Q25 & Q75 & Max \\ 
  \midrule
S0: Non-informative & Intercept (ext) & 200 & 1.99 & 1.95 & 0.28 & 1.29 & 1.80 & 2.16 & 2.92 \\ 
  S0: Non-informative & Poverty (ext) & 200 & 1.90 & 1.89 & 0.32 & 1.23 & 1.67 & 2.09 & 3.07 \\ 
  S0: Non-informative & Urban (ext) & 200 & 3.25 & 3.18 & 0.69 & 1.59 & 2.74 & 3.68 & 5.28 \\ 
  S0: Non-informative & Black (ext) & 200 & 2.08 & 2.05 & 0.36 & 1.14 & 1.83 & 2.29 & 3.46 \\ 
  S0: Non-informative & Hispanic (ext) & 200 & 1.69 & 1.67 & 0.24 & 1.05 & 1.54 & 1.80 & 2.68 \\ 
  S0: Non-informative & Intercept (int) & 200 & 2.00 & 1.97 & 0.30 & 1.38 & 1.79 & 2.17 & 2.92 \\ 
  S0: Non-informative & Poverty (int) & 200 & 1.81 & 1.80 & 0.31 & 1.06 & 1.62 & 1.96 & 3.11 \\ 
  S0: Non-informative & Urban (int) & 200 & 2.93 & 2.86 & 0.70 & 1.41 & 2.44 & 3.33 & 5.35 \\ 
  S0: Non-informative & Black (int) & 200 & 2.07 & 2.02 & 0.39 & 1.32 & 1.77 & 2.28 & 3.55 \\ 
  S0: Non-informative & Hispanic (int) & 200 & 1.61 & 1.59 & 0.23 & 1.11 & 1.45 & 1.76 & 2.34 \\ 
  S0: Non-informative & log kappa & 200 & 1.73 & 1.71 & 0.27 & 1.12 & 1.55 & 1.92 & 2.36 \\ 
  S3: NSECE-calibrated & Intercept (ext) & 200 & 4.02 & 3.88 & 0.97 & 1.97 & 3.29 & 4.71 & 7.47 \\ 
  S3: NSECE-calibrated & Poverty (ext) & 200 & 3.79 & 3.72 & 0.97 & 1.90 & 3.03 & 4.45 & 7.53 \\ 
  S3: NSECE-calibrated & Urban (ext) & 200 & 5.55 & 5.43 & 1.72 & 1.62 & 4.31 & 6.56 & 11.21 \\ 
  S3: NSECE-calibrated & Black (ext) & 200 & 4.01 & 3.81 & 1.09 & 2.07 & 3.13 & 4.64 & 7.23 \\ 
  S3: NSECE-calibrated & Hispanic (ext) & 200 & 3.03 & 2.94 & 0.73 & 1.58 & 2.50 & 3.50 & 5.26 \\ 
  S3: NSECE-calibrated & Intercept (int) & 200 & 3.80 & 3.70 & 0.88 & 1.68 & 3.13 & 4.46 & 6.40 \\ 
  S3: NSECE-calibrated & Poverty (int) & 200 & 3.44 & 3.42 & 0.89 & 1.53 & 2.79 & 3.89 & 7.32 \\ 
  S3: NSECE-calibrated & Urban (int) & 200 & 5.34 & 5.26 & 1.41 & 2.34 & 4.41 & 6.13 & 11.40 \\ 
  S3: NSECE-calibrated & Black (int) & 200 & 3.89 & 3.67 & 1.10 & 1.96 & 3.19 & 4.52 & 7.45 \\ 
  S3: NSECE-calibrated & Hispanic (int) & 200 & 2.72 & 2.65 & 0.63 & 1.67 & 2.20 & 3.09 & 5.16 \\ 
  S3: NSECE-calibrated & log kappa & 200 & 3.12 & 3.07 & 0.69 & 1.38 & 2.74 & 3.55 & 5.44 \\ 
  S4: Stress test & Intercept (ext) & 200 & 5.50 & 5.23 & 1.52 & 2.79 & 4.45 & 6.10 & 13.23 \\ 
  S4: Stress test & Poverty (ext) & 200 & 5.51 & 5.31 & 1.51 & 2.72 & 4.37 & 6.48 & 11.59 \\ 
  S4: Stress test & Urban (ext) & 200 & 6.43 & 6.22 & 1.91 & 2.97 & 5.07 & 7.66 & 17.11 \\ 
  S4: Stress test & Black (ext) & 200 & 5.75 & 5.41 & 1.75 & 2.30 & 4.45 & 6.79 & 11.38 \\ 
  S4: Stress test & Hispanic (ext) & 200 & 4.46 & 4.36 & 1.14 & 1.91 & 3.70 & 5.02 & 10.86 \\ 
  S4: Stress test & Intercept (int) & 200 & 4.86 & 4.71 & 1.26 & 2.38 & 3.93 & 5.52 & 9.53 \\ 
  S4: Stress test & Poverty (int) & 200 & 4.66 & 4.44 & 1.39 & 2.31 & 3.60 & 5.45 & 9.04 \\ 
  S4: Stress test & Urban (int) & 200 & 6.11 & 5.96 & 1.74 & 2.24 & 5.02 & 7.15 & 13.79 \\ 
  S4: Stress test & Black (int) & 200 & 5.19 & 5.01 & 1.48 & 2.13 & 4.16 & 5.98 & 10.13 \\ 
  S4: Stress test & Hispanic (int) & 200 & 3.84 & 3.74 & 0.95 & 1.78 & 3.07 & 4.57 & 6.44 \\ 
  S4: Stress test & log kappa & 200 & 4.21 & 4.10 & 1.06 & 2.11 & 3.50 & 4.85 & 7.69 \\ 
   \bottomrule
\end{tabular}
\caption{Design Effect Ratios (DER) Across Simulation Replications. DER = diag($V_{\text{sand}}$) / diag($H_{\text{obs}}^{-1}$). All 11 fixed-effect parameters shown.} 
\end{table}

}

\paragraph{Cross-margin correlation recovery.}
\cref{tab:rho-cross-sim} reports recovery of the cross-margin
correlation $\rho_{\mathrm{cross}} = \operatorname{cor}(\delta_j^{\mathrm{ext}},
\delta_j^{\mathrm{int}})$ using a plug-in estimator based on posterior
mean random effects across $S = 51$ states.
The unweighted estimator E-UW systematically overestimates
$\rho_{\mathrm{cross}}$, with relative bias growing from $+84\%$ in
S0 to $+127\%$ in S4.
The weighted estimator E-WT shows moderate downward bias ($-26\%$ to
$-45\%$), and E-WS is identical to E-WT because the Cholesky sandwich
correction (\cref{thm:cholesky}) adjusts only fixed-effect covariances.
Per-replication credible intervals are unavailable because full
posterior draws of the correlation matrix $\Omega$ were not retained
in the simulation archive; only frequentist summaries across
replications are reported.

{\small
\begin{table}[t]
\centering
\caption{Recovery of the cross-margin correlation $\rho_{\mathrm{cross}} = \operatorname{cor}(\delta_j^{\mathrm{ext}}, \delta_j^{\mathrm{int}})$
  across $R = 200$ simulation replications. The plug-in estimator 
  $\hat{\rho} = \operatorname{cor}(\bar{\delta}_j^{\mathrm{ext}}, \bar{\delta}_j^{\mathrm{int}})$
  uses posterior mean random effects ($S = 51$ states). True value: $\rho_{\mathrm{cross}} = 0.2853$.
  E\textsubscript{WS} is identical to E\textsubscript{WT} because the sandwich correction applies only to fixed effects.}
\label{tab:rho-cross-sim}
\smallskip
\small
\begin{tabular}{@{}ll rrrrr@{}}
\toprule
Scenario & Estimator & $\overline{\hat{\rho}}$ & Bias & Rel.~Bias (\%) & RMSE & Emp.~SE \\
\midrule
S0 (Baseline)      & E\textsubscript{UW}                 & 0.5238 & $+0.2386$ & $+83.6$ & 0.2990 & 0.1807 \\
                   & E\textsubscript{WT}                 & 0.2104 & $-0.0749$ & $-26.3$ & 0.2294 & 0.2173 \\
                   & E\textsubscript{WS}$^{\dagger}$     & 0.2104 & $-0.0749$ & $-26.3$ & 0.2294 & 0.2173 \\
\midrule
S3 (NSECE)         & E\textsubscript{UW}                 & 0.5814 & $+0.2961$ & $+103.8$ & 0.3643 & 0.2128 \\
                   & E\textsubscript{WT}                 & 0.1658 & $-0.1195$ & $-41.9$ & 0.2256 & 0.1918 \\
                   & E\textsubscript{WS}$^{\dagger}$     & 0.1658 & $-0.1195$ & $-41.9$ & 0.2256 & 0.1918 \\
\midrule
S4 (Heavy-tail)    & E\textsubscript{UW}                 & 0.6468 & $+0.3615$ & $+126.7$ & 0.4075 & 0.1886 \\
                   & E\textsubscript{WT}                 & 0.1579 & $-0.1274$ & $-44.7$ & 0.2372 & 0.2005 \\
                   & E\textsubscript{WS}$^{\dagger}$     & 0.1579 & $-0.1274$ & $-44.7$ & 0.2372 & 0.2005 \\
\bottomrule
\end{tabular}

\medskip
{\footnotesize $^{\dagger}$E\textsubscript{WS} shares the same Stan fit as E\textsubscript{WT}; the Cholesky sandwich correction adjusts only fixed-effect covariances, not hyperparameter posteriors.}
\end{table}

}
